% Paste into examplethesis.tex: replace from \chapter{Method or Methods}
% through the content before \chapter{What you did} (keep \cleardoublepage).
% Locked plan: Code/docs/plan_ch3_ch4.md §7 / §9.
% No result numbers; no success hard-wiring; emotion almost silent.

\chapter{Method or Methods}
\label{ch:methods}

This chapter describes how the degree project investigates the research questions while constructing a research prototype of a booking assistant.
The work follows a design-enquiry approach: an inspectable modular artefact is built, and controlled comparisons are used to gather evidence for RQ1 and RQ2.
The comparisons are exploratory; the method does not presuppose that any modular configuration will outperform a language-model-only baseline.
Section~\ref{sec:method-architecture} presents the system architecture: overall stack, turn-level data flow, and layer-wise design.
Section~\ref{sec:method-data} states the corpus-construction rationale, staged use of public datasets, balancing standards, and reliability checks.
Section~\ref{sec:method-experimental} defines the experimental design.
Section~\ref{sec:method-evaluation} states the evaluation framework, including condition names, metrics, analysis form, and validity limits.

\section{System architecture}
\label{sec:method-architecture}

The architecture separates local intention tracking from surface language generation.
Situational booking goals are maintained in a local sidecar; a large language model produces user-visible replies under a controlled prompt.
Figure~\ref{fig:method-architecture} shows the layered organisation.
On a task turn the utterance is mapped to a structured task variable (STV), an optional aggregate-memory module may supply a prior over bookable goals, a belief tracker updates a multi-turn posterior, dialogue control selects clarification, suggestion, or dual-track observation under topic change, and an intention prompt block assembled from allow-listed fields conditions the language model.
Chitchat is routed separately and does not update belief.
Section~\ref{sec:method-evaluation} defines comparative cuts on intention-block injection and on the memory prior; the layering itself is unchanged across conditions.

\begin{figure}[!ht]
  \centering
  \includegraphics[width=\textwidth]{figures/method_architecture.pdf}
  \caption{Layered architecture of the booking-assistant prototype.
  The local sidecar comprises STV extraction (M1), aggregate-memory priors (M3), belief tracking (M2), dialogue control, and packing of an intention prompt block.
  The language model realises user-visible replies.
  Aggregate profiles and raw identifiers remain outside the generative prompt.
  Comparative cuts are defined in Section~\ref{sec:method-evaluation}.}
  \label{fig:method-architecture}
\end{figure}

M1 maps each utterance to an STV that records a six-class service category (\texttt{fitness}, \texttt{wellness}, \texttt{beauty}, \texttt{medical}, \texttt{dining}, \texttt{travel}), soft constraints such as temporal preference and location, and a category confidence.
The default extractor is a lightweight pipeline model trained on the composite corpus and loaded from a frozen artefact; an LLM-based extractor is retained only as an optional alternative.
Importantly, M1 does not apply an automatic ambiguous gate that rewrites low-confidence predictions into a dialogue-driving \texttt{ambiguous} label.
Underspecified input is therefore left for multi-turn belief updating rather than closed by a single-turn confidence threshold.

When personalisation is enabled, M3 contributes only as a prior over goal hypotheses.
The store retains aggregate counts for category, temporal preference, and price tier, converts them into a prior that is blended before the belief update, and uses a uniform prior when memory is disabled.
Aggregates are written back at session end from the top hypothesis.
User identifiers and aggregate profiles are withheld from the generative prompt, so memory influences generation only through the posterior summarised in the intention block.

M2 represents booking goals as a posterior over a discrete hypothesis space whose factors include service category together with soft constraints such as time, location zone, and price tier.
Each task turn updates this posterior with soft likelihoods dominated by category confidence, optionally strengthened by slot agreement.
Entropy, top-1 mass, a short top-$K$ list, and a category-convergence flag are derived from the posterior and determine whether the surface policy clarifies or suggests.
High residual uncertainty favours clarification; a sufficiently concentrated category mass favours suggestion.
M1 clarification flags are not used as the primary policy trigger.

Dialogue control operates on this belief state under two regimes.
During topic tracking, a clarification protocol can present candidate categories and request confirmation before returning to suggestion.
Under hard topic change, signalled for instance by explicit abandon or pivot phrases, the session opens a dual-track transition window in which the active belief may decay when unsupported while a candidate track accumulates evidence for the new topic, with surface policy typically remaining on observation until promotion, abort, or fallback clarification.
The distinction separates refinement of an incomplete goal from replacement of the goal itself.

After policy resolution, generation is conditioned on an intention prompt block that contains the current message, an STV summary, top hypotheses with weights, a compact belief-state summary, the dialogue policy, and the turn index, optionally together with a short recent-chat window for coherence.
Identifiers, aggregate profiles, full raw STV histories, and the full belief vector remain local.
Disabling intention conditioning yields the LLM-only comparative setting of Section~\ref{sec:method-evaluation}, in which the same generative model runs without that structured block.
In all settings the language model does not update belief; inspectable intention state remains in the sidecar.

\section{Data}
\label{sec:method-data}

Training data must fit the six Fixmeapp booking categories without pretending that public corpora share that taxonomy.
Weather queries, music commands, and clinical symptom talk may look related in a general assistant sense, yet they are not bookable personal services here; including them would add noise rather than signal.
In-domain Fixmeapp-style examples alone are also too few and uneven, especially for dining, travel, fitness, and wellness.
Together with the architectural choice in Section~\ref{sec:method-architecture} to resolve underspecification in multi-turn belief, this implies a clear plan: keep \texttt{ambiguous} out of training labels; start from domain seeds and only the public intents that map cleanly; then widen coverage where the mapping remains defensible; finally balance the mix and check it with resources that separate remapped public text from deliberately vague in-domain hard cases.
Row counts, scripts, and full mapping tables are deferred to Chapter~\ref{ch:whatYouDid}.

The first public layer privileges precision over coverage.
ATIS~\cite{hemphill-etal-1990-atis} supplies travel queries, remapped as \texttt{travel} and retained as an external check for that class.
SNIPS~\cite{coucke_snips_2018} contributes only restaurant-booking utterances to \texttt{dining}; its weather, music, and similar intents are skipped.
Remaining classes lean on in-domain seeds until a second layer fills the gaps.
Schema-Guided Dialogue~\cite{rastogi_sgd_2020} adds salon- and clinic-like service turns when the service name maps stably to beauty or medical, while media and messaging domains are ignored.
The Hybrid Appointment Scheduling Dataset~\cite{hasd_hf} contributes healthcare schedule and cancel intents, not open medical chat.
LUCID~\cite{stacey_lucid_2024} adds fitness and wellness booking-style turns extracted from user queries paired with system intents.
If a class is still thin after deduplication, short template expansions from in-domain phrasing are allowed and labelled as such, so they never enter the hard-case evaluation set unnoticed.
Symptom-triage corpora are not used for \texttt{medical}, because diagnosis dialogue is the wrong task for booking extraction.

After merging, duplicates are removed and each class is capped (default four hundred examples), with in-domain rows preferred over public fills.
The cap and source priority follow from an early dining-heavy SNIPS merge in which one abundant proxy dominated the class prior.
Fixed-seed splits keep the build reproducible.
Reliability is assessed rather than assumed: ATIS travel performance tests the travel mapping; held-out in-domain and hard-case slices test Fixmeapp-like and underspecified language; composite holdouts, when reported, describe the remapped mix and must not be interpreted as field performance on Fixmeapp users.
These checks establish whether the structured front-end is usable as evidence for later stages; they do not answer RQ1 or RQ2.

Dialogue-level evaluation does not use live Fixmeapp conversations.
Scripted multi-turn scenarios instead cover the six categories and include vague openings, constraint revision, and topic change, with a session-level gold category hint.
Vague input for RQ1 and RQ2 is therefore operationalised by controlled scenario design, not by assigning a vague label to every utterance.
Public corpora are used under their research licences; real customer booking histories are not collected.

\section{Experimental design}
\label{sec:method-experimental}

Evidence is collected at three levels: component checks, multi-turn scenario replay, and condition contrasts for RQ1 and RQ2.
The comparisons are exploratory; modular settings are not assumed to outperform the language-model-only baseline.
Metrics and validity limits are defined in Section~\ref{sec:method-evaluation}.

Component checks evaluate STV extraction on held-out in-domain data, hard cases, and mapped public subsets such as ATIS travel.
They establish whether the structured front-end yields usable evidence for later stages and do not answer RQ1 or RQ2.

Multi-turn scenario replay runs the scripted sessions of Section~\ref{sec:method-data} through belief updating and dialogue control.
The observables are whether uncertainty narrows, when clarification is preferred over suggestion, and whether a session meets the resolution criterion in Section~\ref{sec:method-evaluation}.
A secondary \textbf{STV-only} policy baseline (per-turn structured evidence without full multi-turn belief policy) may be reported for diagnosis; it is not a thesis-facing language-model condition.

Condition contrasts reuse the same evaluation items under the settings in Table~\ref{tab:method-conditions}.
Prose uses the thesis names; code identifiers are retained for reproducibility.
RQ1 contrasts \textbf{LLM-only} with \textbf{Intention-conditioned}.
RQ2 contrasts \textbf{Intention-conditioned} with \textbf{Memory-augmented}, where an aggregate prior is applied to belief while identifiers and aggregate profiles remain outside the generative prompt.
A null or negative behavioural difference remains informative if the privacy constraint holds.

\begin{table}[!ht]
  \centering
  \small
  \caption{Comparative conditions (thesis names and code identifiers).}
  \label{tab:method-conditions}
  \setlength{\tabcolsep}{4pt}
  \begin{tabular}{@{}lp{0.28\textwidth}p{0.42\textwidth}@{}}
    \hline
    \textbf{Thesis name} & \textbf{Code id} & \textbf{Difference} \\
    \hline
    LLM-only & \texttt{C0\_llm\_only} & No intention summary in the prompt \\
    Intention-conditioned & \texttt{C1\_intention\_block} & Intention summary in the prompt \\
    Memory-augmented & \texttt{C2\_m3\_prior} & As above, plus aggregate prior on belief \\
    \hline
  \end{tabular}
\end{table}

Runs use offline evaluation scripts and, when needed, an interactive prototype for qualitative inspection.
Software entry points are those of the project repository; hardware is a standard development workstation.
Decoding randomness is treated as a validity threat in Section~\ref{sec:method-evaluation}.

\section{Evaluation framework}
\label{sec:method-evaluation}

The evaluation criteria are chosen by reference to prior task-oriented and clarification work: where a standard fits this prototype it is kept; where it does not fit, it is dropped or replaced; where the research questions require an extra check, it is added.
Table~\ref{tab:method-metrics} makes those decisions explicit.

Task-oriented evaluation in MultiWOZ asks whether the system reaches the user goal and how many turns it needs~\cite{budzianowski_multiwoz_2018}.
That concern is kept here as category-level task match on shared items, which is the primary outcome for RQ1 and RQ2, together with turns to resolution as an efficiency process measure.
Belief-state dialogue work treats posterior uncertainty as central to policy~\cite{young_pomdp_2013}, so entropy is kept to observe whether belief narrows under vague sessions.
Clarification studies score the decision to ask rather than answer~\cite{aliannejadi_clarifying_2019,lin_ask-to-clarify_2025}, so clarify-versus-suggest rate is kept under underspecification.
Intent accuracy from ATIS and SNIPS~\cite{hemphill-etal-1990-atis,coucke_snips_2018} is kept only as a front-end check and does not answer the research questions.

Two common standards are not used.
MultiWOZ Inform and Success assume API or database ground truth for booking completion~\cite{budzianowski_multiwoz_2018}; this prototype has no such backend, so those scores are replaced by category match rather than reported as Inform/Success.
Human preference judgements, frequent in recent LLM dialogue papers, are also omitted because the study prioritises inspectable structural contrasts over subjective ranking of surface text.

One check is added beyond standard TOD benchmarks.
Those benchmarks do not ask whether aggregate memory is withheld from the generative model, yet RQ2 requires that constraint.
A prompt privacy audit is therefore included, verifying that raw identifiers and serialised detailed profiles are absent from the intention block and generative input~\cite{staufer_name_audit_2026}.
For RQ2 this audit is primary; a large personalisation gain is not required if the privacy constraint holds.

\begin{table}[!ht]
  \centering
  \small
  \caption{Evaluation criteria relative to established standards.}
  \label{tab:method-metrics}
  \setlength{\tabcolsep}{3pt}
  \begin{tabular}{@{}p{0.26\textwidth}ccp{0.38\textwidth}@{}}
    \hline
    \textbf{Criterion} & \textbf{In prior work} & \textbf{In this study} & \textbf{Rationale} \\
    \hline
    Goal / task attainment & Yes & Retain & Primary RQ1/RQ2 outcome as category match~\cite{budzianowski_multiwoz_2018} \\
    Turns / efficiency & Yes & Retain & Process measure: turns to resolution~\cite{budzianowski_multiwoz_2018} \\
    Belief uncertainty & Yes & Retain & Process measure: posterior entropy~\cite{young_pomdp_2013} \\
    Ask vs.\ answer & Yes & Retain & Process measure: clarify vs.\ suggest~\cite{aliannejadi_clarifying_2019} \\
    Intent accuracy & Yes & Retain$^*$ & Front-end only~\cite{hemphill-etal-1990-atis,coucke_snips_2018} \\
    Inform / Success & Yes & Omit & No booking API or DB ground truth \\
    Human preference & Yes & Omit & Structural contrasts preferred \\
    Prompt privacy audit & No & Introduce & Required by RQ2~\cite{staufer_name_audit_2026} \\
    \hline
  \end{tabular}\\[0.3em]
  {\footnotesize $^*$Background evidence; not a substitute for condition contrasts.}
\end{table}

RQ1 contrasts \textbf{LLM-only} with \textbf{Intention-conditioned} on the kept behavioural quantities.
RQ2 contrasts \textbf{Intention-conditioned} with \textbf{Memory-augmented} on the same quantities plus the privacy audit.
Results appear as condition contrasts with brief notes on unstable cases.
Reproducibility rests on fixed scripts, condition definitions, and versioned data.
Limits include scripted scenarios, decoding variability, and public-corpus domain shift; STV-only diagnostics must not be read as LLM-only.
Absence of improvement under modular settings remains a valid outcome.

\cleardoublepage
